% Appendices for the final SmogShield report.

\section{Viva questions and answers}
\label{app:viva}

\begin{description}[leftmargin=1.2em,style=nextline,itemsep=4pt]
  \item[What does the app predict?]
  It predicts tomorrow's PM$_{2.5}$ concentration and tomorrow's AQI category
  from six values for today: date, PM$_{2.5}$, AQI category, temperature,
  humidity, and wind speed. The forecast is always for the day after the
  selected feature date.

  \item[Why are there two AQI category views?]
  The classifier is a direct six-class model. The implied category converts the
  regression's predicted PM$_{2.5}$ through the same US-EPA breakpoints used to
  build the labels, so it is monotonic in the prediction. The classifier was
  weaker than persistence on test, so the app treats the implied category as the
  more defensible view and warns when the two disagree.

  \item[Does the model beat persistence?]
  No. On test, regression RMSE is 64.92 versus 64.82 for persistence, a
  0.10 $\mu$g/m$^3$ difference and a skill score of $-0.0016$. Classification
  macro-F1 is 0.4780 versus 0.5157. Ridge reached a lower test RMSE (62.10), but
  it was not the validation-selected model, so switching to it after seeing test
  would violate the validation-only rule.

  \item[Why keep a model that does not beat persistence?]
  The project is an end-to-end explainable system, not an accuracy record. The
  model was selected on validation, the test result is reported honestly, and
  the application demonstrates inference, uncertainty, explanations, and
  reproducibility. Persistence is a strong one-day baseline because PM$_{2.5}$ is
  highly autocorrelated; failing to beat it is an informative result.

  \item[What is HGB?]
  Histogram-based gradient boosting. It builds an ensemble of small decision
  trees sequentially, with each tree correcting the errors of the previous ones.
  The histogram step bins continuous features for speed. It handles non-linear
  relationships, interactions, and missing values, and it needs little
  preprocessing. The frozen settings were selected on validation in Deliverable
  2.

  \item[What is persistence?]
  Persistence assumes tomorrow will be the same as today: predicted value at
  $t+1$ equals the observed value at $t$. For classification, tomorrow's
  category equals today's. It is the simplest possible forecast and the baseline
  every model must beat.

  \item[Why macro-F1 instead of accuracy?]
  Accuracy is dominated by the common \emph{Unhealthy} class. Macro-F1 gives each
  class equal weight, so it exposes poor performance on the rare hazardous
  classes. The HGB classifier has slightly higher test accuracy than persistence
  but lower macro-F1; reporting only accuracy would obscure the rare-class
  weakness.

  \item[What does the 80\% interval mean?]
  It is a central interval for the next-day PM$_{2.5}$ observation, built from
  quantile HGB models at the 10th and 90th percentiles and calibrated on
  validation. Validation coverage is 80.3\%; test coverage is 77.0\%. The test
  split is used only to report coverage, never to choose the calibration factor.

  \item[Why is the interval so wide?]
  The model is uncertain, especially on severe days. The default November
  example forecasts 189.74 $\mu$g/m$^3$ with an interval of 133.8--294.1. The
  width is an honest statement about the evidence, not a bug. The backtest on
  1 November 2024 shows that the actual value can fall outside the interval when
  the day is an extreme outlier.

  \item[Why only six inputs?]
  Usability. The remaining features use training medians or modes, and the
  PM$_{2.5}$ lags and rolling windows are held at today's value. The app calls
  itself an illustrative single-day forecaster for this reason. A production
  version would ingest live weather and a multi-day PM$_{2.5}$ history.

  \item[What changed after Deliverable 2?]
  The model and its reported scores are unchanged; the D3 refit reproduces D2
  exactly. What is new is the application layer: inference, the Streamlit app,
  prediction intervals, explanations, the ablation, the hourly profile, the
  seasonal defaults, and the tests. The D2 code and report were not edited.

  \item[What is the biggest limitation?]
  The single monitoring station and the severe-tail under-prediction. The model
  is weakest on the highest-pollution days, which are exactly the days that
  matter most for public health. A second station, live inputs, and a
  cost-sensitive or conformal approach are the next steps.
\end{description}

\section{Glossary}
\label{app:glossary}

\begin{longtable}{@{}L{3.6cm}L{11.0cm}@{}}
\toprule
\textbf{Term} & \textbf{Meaning in this project} \\
\midrule
\endhead
\bottomrule
\endfoot
PM$_{2.5}$ & Fine particulate matter with diameter at most 2.5 $\mu$m, measured
in $\mu$g/m$^3$. The regression target. \\
AQI & Air Quality Index, a 0--500 scale computed from PM$_{2.5}$ using the US-EPA
2012 breakpoints. \\
HGB & Histogram-based gradient boosting; an ensemble of small trees trained
sequentially to correct errors. \\
Regression & Predicting a continuous number, here tomorrow's PM$_{2.5}$. \\
Classification & Predicting a discrete class, here tomorrow's AQI category. \\
Persistence & The baseline ``tomorrow equals today.'' \\
RMSE & Root mean squared error; same units as the target and sensitive to large
errors. \\
MAE & Mean absolute error; average absolute miss, in the target's units. \\
$R^2$ & Fraction of variance explained; 1 is perfect and 0 is no better than the
mean. \\
Macro-F1 & The unweighted mean of per-class F1 scores; treats each class equally
and exposes rare-class performance. \\
Weighted-F1 & F1 weighted by class support; dominated by the common classes. \\
Precision and recall & Of the predicted positives, how many are correct; of the
true positives, how many were found. \\
Confusion matrix & A table of predicted versus true classes. \\
Prediction interval & A range for a future observation, here the central 80\%. \\
Quantile regression & Training a model to predict a conditional percentile rather
than the mean. \\
Coverage & The fraction of true values that fall inside the interval. \\
Calibration & Adjusting interval width so empirical coverage matches the nominal
level. \\
Permutation importance & Global importance: the increase in error when a feature
is shuffled. \\
Sensitivity & Local importance: the change in the forecast when one input is
perturbed. \\
Ablation & Removing feature groups to test whether they improve performance. \\
Data leakage & Information from validation or test contaminating training;
prevented here by the seven-day purge and train-only preprocessing. \\
Class imbalance & Some categories are much rarer than others; handled with
balanced class weights and macro-F1. \\
Piecewise constant & A tree model's prediction stays flat inside a leaf, so a
small input change may have no effect. \\
Steady state & The app's assumption that today's PM$_{2.5}$ has been stable for
the past week when it fills lags and rolling windows. \\
\end{longtable}

\section{File map and build notes}
\label{app:files}

Table~\ref{tab:filemap} lists the Deliverable 3 files a reviewer is most likely
to open. The full repository map also includes the Deliverable 1 data scripts and
the Deliverable 2 modelling code, which are unchanged.

\begin{table}[htbp]
\centering
\footnotesize
\begin{tabular}{@{}L{5.4cm}L{9.2cm}@{}}
\toprule
\textbf{File} & \textbf{Role} \\
\midrule
\texttt{app/app.py} & Streamlit interface: six inputs, forecast cards, interval,
guidance, charts, expanders \\
\texttt{src/inference.py} & Loads the four pipelines and metadata; assembles the
feature row; predicts PM$_{2.5}$, category, interval, explanation, and guidance \\
\texttt{src/train\_pipeline.py} & Refits the frozen pipelines; writes the joblib
artifacts and \texttt{metadata.json} \\
\texttt{src/ablation.py} & Runs the nested history/calendar/weather ablation \\
\texttt{src/hourly.py} & Builds the hour-of-day and seasonal PM$_{2.5}$ profile \\
\texttt{tests/test\_inference.py} & Fourteen unit tests for the inference and app
contract \\
\texttt{report/make\_figures.py} & Generates the report figures and the worked
examples from the frozen artifacts \\
\texttt{report/report\_d3.tex} & This report's master LaTeX file \\
\texttt{report/phase34\_final.tex} & Deliverable 2 Phase 3--4 body, reused with
adjusted table paths \\
\texttt{report/phase56\_final.tex} & New Deliverable 3 Phase 5--6 chapters \\
\texttt{report/appendix\_final.tex} & Viva questions, glossary, and file map \\
\texttt{README.md} & Install, run, artifact, and submission notes \\
\texttt{STUDY\_REPORT.md} & Plain-language technical study report \\
\texttt{requirements.txt} & Python dependencies for the application layer \\
\bottomrule
\end{tabular}
\caption{Deliverable 3 file map.}
\label{tab:filemap}
\end{table}

\paragraph{Build notes.} The report is compiled from the
\texttt{deliverable\_3/report} directory with Tectonic:
\begin{verbatim}
tectonic report_d3.tex
\end{verbatim}
The \texttt{report/make\_figures.py} script must be run first if the figures or
\texttt{examples.json} are missing. The model artifacts must be present for the
script to compute the worked examples; run \texttt{python -m src.train\_pipeline}
from \texttt{deliverable\_3} first. The report's numerical tables are taken from
the frozen D2 metrics and the D3 metadata, ablation, and worked-example files.
